\setcounter{ex}{0}

% --- CÂU 1 ---
\begin{ex}
	Gọi $M, m$ lần lượt là giá trị lớn nhất, giá trị nhỏ nhất của hàm số $y = \cos^2 x + 2\sin x$. Tính $M - m$.
	\shortans{4}
	\loigiai{
		Ta có:
		$$y = \cos^2 x + 2\sin x = 1 - \sin^2 x + 2\sin x = 2 - (\sin x - 1)^2.$$
		Với mọi $x \in \mathbb{R}$, vì $-1 \le \sin x \le 1$ nên $-2 \le \sin x - 1 \le 0 \implies 0 \le (\sin x - 1)^2 \le 4$.\\
		Suy ra:
		$$-2 \le 2 - (\sin x - 1)^2 \le 2 \implies -2 \le y \le 2.$$
		Do đó $M = 2$ (đạt được khi $\sin x = 1$) và $m = -2$ (đạt được khi $\sin x = -1$).\\
		Vậy $M - m = 2 - (-2) = 4$.
	}
\end{ex}

% --- CÂU 2 ---
\begin{ex}
	Cho $\cos\alpha = -\dfrac{3}{5}$ và $\pi < \alpha < \dfrac{3\pi}{2}$. Tính giá trị của biểu thức $B = 5\sin\alpha - 4\cot\alpha$.
	\shortans{-7}
	\loigiai{
		Vì $\pi < \alpha < \dfrac{3\pi}{2}$ nên $\sin\alpha < 0$. Ta có:
		$$\sin\alpha = -\sqrt{1 - \cos^2\alpha} = -\sqrt{1 - \left(-\dfrac{3}{5}\right)^2} = -\dfrac{4}{5}.$$
		Suy ra $\cot\alpha = \dfrac{\cos\alpha}{\sin\alpha} = \dfrac{-\frac{3}{5}}{-\frac{4}{5}} = \dfrac{3}{4}$.\\
		Thay vào biểu thức $B$, ta được:
		$$B = 5 \cdot \left(-\dfrac{4}{5}\right) - 4 \cdot \dfrac{3}{4} = -4 - 3 = -7.$$
	}
\end{ex}

% --- CÂU 3 ---
\begin{ex}
	Một rạp hát có $20$ hàng ghế, hàng thứ nhất có $18$ ghế, mỗi hàng sau nhiều hơn hàng liền trước $2$ ghế. Rạp hát có tất cả bao nhiêu ghế?
	\shortans{740}
	\loigiai{
		Số ghế ở mỗi hàng từ hàng thứ nhất đến hàng thứ $20$ lập thành một cấp số cộng có số hạng đầu $u_1 = 18$, công sai $d = 2$ và số các số hạng là $n = 20$.\\
		Tổng số ghế trong rạp hát là:
		$$S_{20} = 20u_1 + \dfrac{20 \cdot 19}{2}d = 20 \cdot 18 + 190 \cdot 2 = 360 + 380 = 740\text{ (ghế)}.$$
	}
\end{ex}

% --- CÂU 4 ---
\begin{ex}
	Chiều cao mực nước biển (mét) tại thời điểm $t$ (giờ) trong một ngày được mô hình hóa bởi hàm số $h(t) = 3 + 2\sin\left(\dfrac{\pi}{6}t\right)$. Tìm thời điểm đầu tiên $t > 0$ để chiều cao mực nước đạt $4\text{ m}$.
	\shortans{1}
	\loigiai{
		Để chiều cao mực nước đạt $4\text{ m}$ thì:
		$$h(t) = 4 \iff 3 + 2\sin\left(\dfrac{\pi}{6}t\right) = 4 \iff \sin\left(\dfrac{\pi}{6}t\right) = \dfrac{1}{2}$$
		$$\iff \hoac{\dfrac{\pi}{6}t &= \dfrac{\pi}{6} + k2\pi \\ \dfrac{\pi}{6}t &= \dfrac{5\pi}{6} + k2\pi} \iff \hoac{t &= 1 + 12k \\ t &= 5 + 12k}\ (k \in \mathbb{Z}).$$
		Thời điểm đầu tiên $t > 0$ ứng với $k = 0$ ở nhánh thứ nhất là $t = 1$ (giờ).
	}
\end{ex}

% --- CÂU 5 ---
\begin{ex}
	Cho hình chóp $S.ABCD$ có đáy $ABCD$ là hình bình hành tâm $O$, $P$ là trung điểm của $SD$, $J = BP \cap SO$. Biết $\dfrac{SJ}{SO} = \dfrac{a}{b}$ với $\dfrac{a}{b}$ là phân số tối giản. Tính $a + b$.
	\shortans{5}
	\loigiai{
		Vì $ABCD$ là hình bình hành tâm $O$ nên $O$ là trung điểm của $BD$.\\
		Xét tam giác $SBD$ có hai đường trung tuyến $SO$ và $BP$ cắt nhau tại $J$, suy ra $J$ là trọng tâm của tam giác $SBD$.\\
		Theo tính chất trọng tâm tam giác, ta có:
		$$\dfrac{SJ}{SO} = \dfrac{2}{3}.$$
		Do $\dfrac{2}{3}$ là phân số tối giản nên $a = 2, b = 3 \implies a + b = 2 + 3 = 5$.
	}
\end{ex}

% --- CÂU 6 ---
\begin{ex}
	Một vật dao động điều hòa có li độ (cm) tại thời điểm $t$ (giây) là $x = 6\sin t + 8\cos t$. Tính biên độ dao động của vật (đơn vị: cm).
	\shortans{10}
	\loigiai{
		Ta có:
		$$x = 6\sin t + 8\cos t = \sqrt{6^2 + 8^2}\left(\dfrac{6}{10}\sin t + \dfrac{8}{10}\cos t\right) = 10\cos(t - \varphi),$$
		trong đó $\cos\varphi = \dfrac{8}{10} = \dfrac{4}{5}$ và $\sin\varphi = \dfrac{6}{10} = \dfrac{3}{5}$.\\
		Vậy biên độ dao động của vật là $A = \sqrt{6^2 + 8^2} = 10\text{ (cm)}$.
	}
\end{ex}
